\begin{figure*}[t]
\centering
\begin{tcolorbox}[
  colback=gray!5,colframe=gray!50,boxrule=0.4pt,
  fontupper=\small\ttfamily,left=6pt,right=6pt
]
You are answering a question about a long document. You cannot see the
whole document. EVIDENCE is arranged one document page at a time, in
document order. A block may contain a READER'S RECORD, PARSED TEXT, and a
PAGE IMAGE. Reader records may be incomplete or mistaken. Use the parsed
text to check words, names, and numbers, and the image to check visual
and spatial details.\\[0.5em]

Use only the EVIDENCE. Think step by step, write that thinking after
REASONING:, and give a concise final answer after ANSWER:.\\[0.5em]

Check the question's premises before answering: verify the particular
entities, dates, locations, figures, groups, and measurements it assumes.
Do not substitute a nearby fact when a required detail conflicts with
the evidence. Combine relevant facts across page blocks. For a count
or list, identify matching items page by page, exclude duplicates and
irrelevant items, and then aggregate them.\\[0.5em]

Answer when the evidence supports an answer. Write ``Not answerable''
only when a required fact is genuinely absent or contradicted. Before
doing so, inspect the available parsed text and images and write:
MISSING: <the specific unsupported fact>.\\[0.5em]

REASONING: <premise checks and evidence integration>\\
ANSWER: <concise answer or exactly ``Not answerable''>\\[0.5em]

QUESTION: \{question\}\\
EVIDENCE: \{page blocks containing records and, for flagged pages,
parser text and images\}
\end{tcolorbox}
\caption{Synthesizer prompt structure for iteration 7. If the reader flags no page, the evidence includes the two highest-ranked inspected pages in full.}
\label{fig:app_synthesizer_prompt}
\end{figure*}